%% ===================================================================
%%  SigVault -- tool paper, JOSS-like layout.
%%
%%  Layout lives in jossalike.sty; this file is content only.
%%  SKELETON / GUIDE DRAFT. Not submission copy.
%%
%%  Every \guide{...} block is scaffolding addressed to the author.
%%  Set \showguidesfalse to hide them all and read the prose alone.
%%
%%  Build:  latexmk -pdf paper.tex
%% ===================================================================

\documentclass[10pt,a4paper]{article}
\usepackage{jossalike}

% ---- author-facing guidance ----------------------------------------
\definecolor{guidecol}{HTML}{8A5A00}
\newif\ifshowguides
\showguidestrue                 % <-- \showguidesfalse to hide all guide notes
\newcommand{\guide}[1]{%
  \ifshowguides
    {\par\smallskip\footnotesize\color{guidecol}%
     \textbf{[guide]}~#1\par\smallskip}%
  \fi}
\newcommand{\needcite}{\textcolor{red}{[CITE]}}
\newcommand{\needdata}{\textcolor{red}{[DATA]}}

% Stand-in for a figure whose artwork does not exist yet.
\newcommand{\figstub}[1]{%
  \fbox{\parbox[c][1.55in][c]{\dimexpr\textwidth-2\fboxsep-2\fboxrule\relax}%
    {\centering\color{red}\footnotesize [FIGURE] #1}}}

% ====================================================================
%  MASTHEAD AND METADATA
%  Every field below is a placeholder. Fill in the ones that become
%  real; delete the sidebar blocks you will never have.
% ====================================================================
\renewcommand{\journalmasthead}{%
  {\fontsize{30}{32}\selectfont\bfseries \japlaceholder{VENUE}}\\[3pt]
  {\footnotesize\japlaceholder{Replace with the venue wordmark, or
   \textbackslash includegraphics\{logo.pdf\}.}}}

\renewcommand{\paperdoi}{\japlaceholder{10.xxxxx/xxxxxx}}
\renewcommand{\paperdoiurl}{}
\renewcommand{\softwarerepo}{https://github.com/\japlaceholder{org/sigvault}}
\renewcommand{\softwarereview}{}
\renewcommand{\softwarearchive}{}
\renewcommand{\handlingeditor}{\japlaceholder{unassigned}}
\renewcommand{\reviewerlist}{}
\renewcommand{\datesubmitted}{\japlaceholder{DD Month 2026}}
\renewcommand{\datepublished}{\japlaceholder{DD Month 2026}}

\renewcommand{\jafootline}{Redmond \& Dasgupta, (2026). SigVault: An
  Interactive Relational Taxonomy of RF Signal Security in Space
  Communications.}
\renewcommand{\jafootsource}{\japlaceholder{\textit{Venue}, vol(iss), id.}}

\renewcommand{\papertitle}{SigVault: An Interactive Relational Taxonomy
  of RF Signal Security in Space Communications}

\renewcommand{\paperauthors}{Nicholas D. Redmond\textsuperscript{1},
  Dipankar Dasgupta\textsuperscript{1}}

\renewcommand{\paperaffiliations}{%
  \textbf{1} Department of Computer Science, University of Memphis,
  Memphis, TN, USA}

\begin{document}
\maketitle

% ====================================================================
\section{Summary}

The literature on radio-frequency (RF) security is fragmented across
terrestrial and orbital domains and often confined to the specific time
it was written, moving slower than the attack surface is expanding. With
AI pushing the frontiers of signal exploitation and serious hardware
becoming cheaper and more accessible, the need for a modern and live
taxonomy of RF Signal Security is obvious. We present SigVault, an
interactive relational taxonomy in which knowledge of RF security becomes
more accessible and engaging than ever before. Featured is a matrix of
eleven signal-security facets as they relate to one another, rendered on
a rotatable surface intended to map the literature in a way that is both
navigable and visually pleasing. Additionally, three coupled drill-down
views are provided to allow users to visualize and interact with the
details of the taxonomy in a memorable way. These views are named for
their domain and purpose: TRANSEC, COMSEC, and our Signal Lab. This
taxonomy is designed to reference the literature as well as map the
research gaps in the field, shortening the time it takes to even find
open problems. It is intended to be a living document that will continue
to grow with the field, which is essential in a time where AI advances
are so rapid that the domain has changed before the academic journals can
review and publish the findings.

\guide{Trim to roughly 150--200 words before submission; the version
above is long on purpose so you can cut rather than invent. The sentence
that must survive intact is the one naming the contribution:
\emph{crossing} facets that prior work treats separately. Three edits
already applied from the IEEE draft: \emph{rotatble}/\emph{navigatable}/%
\emph{essentialy} spellings, and ``our proprietary Signal Lab'' ---
proprietary means closed and commercial, which contradicts the living,
citable, open instrument the rest of this paragraph argues for.}

% ====================================================================
\section{Statement of Need}

The security of a radio link is decided by properties that no single
research community owns. Whether a transmission can be detected depends
on its waveform and its antenna pattern; whether it can be exploited
depends on its framing and its encryption; whether it can be denied
depends on the geometry between transmitter, receiver, and adversary.
These properties are studied in electronic warfare, in satellite systems
engineering, in physical-layer security, and in mission cyber defence,
and each community has its own vocabulary, its own canonical threats, and
its own reference documents.

\guide{This is the Introduction from the IEEE draft, carried over intact.
In the JOSS idiom this section has one job: establish that a practitioner
is currently blocked, and that no existing artifact unblocks them. The
paragraph above sets up the fragmentation; it still needs the second
half --- \emph{what a researcher cannot do today}, stated as a concrete
failed task, and why the existing surveys and knowledge bases do not
answer it.}

\begin{figure}[t]
  \centering
  \figstub{Home view --- the eleven-facet N\textsuperscript{2} matrix
    wrapped on the rotatable cylinder, with the crosshair and pair
    readout active.}
  \caption{\japlaceholder{Caption pending.} The Home view. Both axes
    carry the same eleven facets, so every cell is a relation between two
    facet types; column order follows the countermeasure loop, which is
    what makes the axis cyclic and the surface closed.}
  \label{fig:home}
\end{figure}

\guide{A tool paper in this format is carried by its figures --- the
reference layout runs four across eight pages, each one load-bearing
rather than decorative. Budget at least three: the Home cylinder
(Figure~\ref{fig:home}), a crossing modal showing what a cell actually
yields, and one drill-down view. Caption them so a reader who only
looks at figures still gets the argument.}

% ====================================================================
\section{Related Work}

\subsection{Structured Adversary Knowledge Bases}

SPARTA \needcite{}

\guide{\citet{remy2026systematictaxonomyattacksspace} is your closest
neighbour and a reviewer will reach it before you name it. Distinguish
explicitly and early: that work is a systematic taxonomy \emph{document}
at the infrastructure level; this is a queryable relational
\emph{instrument} at the physical layer. Say so in this subsection, not
in Limitations.}

\subsection{Surveys and Taxonomies of Satellite and PHY Security}

Two limitations recur. First, these treatments are largely
orbit-agnostic, which understates how strongly the threat model depends
on regime. Second, they are organised as prose surveys, so the knowledge
they contain cannot be queried along an axis the author did not
anticipate.

\subsection{Matrix and Adjacency Visualization}

\guide{Cite the DSM/N\textsuperscript{2} chart convention and any prior
work laying a matrix on a curved or 3D surface. If no such prior work
exists, say that plainly --- it is a stronger sentence than a vague
gesture at ``related visualization work''.}

% ====================================================================
\section{Taxonomy Structure}

\subsection{Facets}

\begin{table}[t]
\caption{The eleven facets and their families. Family membership drives
hue in the visual encoding.}
\label{tab:facets}
\centering
\footnotesize
\begin{tabular}{@{}lll@{}}
\toprule
\textbf{Facet} & \textbf{Family} & \textbf{Example members} \\
\midrule
Protocol     & System     & DVB-S2X, Starlink Ku, GPS L1 \\
Use-case     & System     & broadcast, TT\&C, backhaul \\
Band         & Spectrum   & HF, L, S, X, Ku/Ka, mmWave, optical \\
Modulation   & Spectrum   & BPSK, OQPSK, BOC(10,5), DSSS \\
Antenna      & Physical   & helical, parabolic, phased array \\
Propagation  & Physical   & vacuum, troposphere, multipath \\
Noise        & Physical   & thermal, co-channel, intentional \\
Attack       & Adversary  & intercept, follower jamming, meaconing \\
Objective    & Defence    & LPD, LPI, LPE, AJ, TFS \\
Mitigation   & Defence    & spreading, nulling, OSNMA \\
Limitation   & Residual   & PAPR cost, acquisition delay \\
\bottomrule
\end{tabular}
\end{table}

\subsection{Relations as the Unit of Knowledge}

\subsection{Evidence Density}

\guide{\textbf{Load-bearing, and currently unanswered.} Evidence density
is the quantity the entire visual encoding maps to luminance, so ``how
did you count?'' is the first question a reviewer asks. Three
operationalisations are available: raw citation count per cell;
normalised count against facet-pair prior; or an ordinal rubric
(none~/~anecdotal~/~empirical~/~replicated). Recommend the ordinal
rubric with two independent raters and a reported Cohen's~$\kappa$ ---
it is the only option that survives the question ``why is this cell
darker than that one?''}

% ====================================================================
\section{Visual Design}
\label{sec:vis}

\subsection{Requirements}

\subsection{A Cyclic Axis}

\guide{This subsection carries the cylinder. The claim is that column
order follows the countermeasure loop --- Protocol $\to$ Band $\to$
Modulation $\to$ Antenna $\to$ Propagation $\to$ Noise $\to$ Attack
$\to$ Objective $\to$ Mitigation $\to$ Limitation $\to$ Use-case $\to$
Protocol --- so the axis genuinely has no first element and a closed
surface is the honest projection. If that ordering cannot be defended,
the honest fallback is a flat matrix, and the paper is built so losing
this argument costs one subsection rather than the contribution. Keep
that property as you write.}

\subsection{Encoding}

\subsubsection{Rejected: identity as a colour blend}

\subsection{Interaction}

% ====================================================================
\section{System Architecture}

\subsection{One Graph, Four Projections}

\subsection{Provenance and Reproducibility}

\guide{State the snapshot-not-fetch policy here: external sources are
version-pinned into the repository with a retrieval date rather than
queried live, so a figure in this paper can be regenerated years from
now. Name the snapshot commit. This is the subsection that earns the
word ``living'' used in the Summary --- without it, the claim is
decoration.}

\subsection{Implementation}

% ====================================================================
\section{Usage Scenarios}

\guide{Two scenarios is the right number at this length; three if you cut
elsewhere. Each should be a task a real analyst has, not a feature tour,
and each should end with an answer the reader could not have obtained
from the sources cited in Related Work.}

\subsection{Locating a Research Gap}

\subsection{GNSS Spoofing}

% ====================================================================
\section{Evaluation Plan}

\subsection{Classroom Deployment}

\subsection{Expert Walkthrough}

\subsection{Task-Based Study}

\guide{\textbf{Load-bearing, and currently unrun.} There is no
cylinder-versus-flat-grid comparison anywhere in this work. Until one
exists, the curved surface is an aesthetic assertion. The minimum
credible study is a within-subjects lookup task --- ``find the cell for
facet~A~$\times$~facet~B'' --- timed on both layouts, and you should be
willing to report the result even if the grid wins.}

\subsection{Case-Study Validation}

% ====================================================================
\section{Limitations}

% ====================================================================
\section{Conclusion}

% ====================================================================
\section{Acknowledgements}

The authors thank \needdata{}.

% ====================================================================
\nocite{*}   % remove once \needcite{} markers become real \cite{} calls
\bibliographystyle{plainnat}
\bibliography{refs}

\end{document}
